%------------------------------------------------------------------------------%
\documentclass[a4paper,12pt,fleqn]{article}

\usepackage{gaal}

\ShowAnswers

%------------------------------------------------------------------------------%
\begin{document}

\makeHeader{GAAL}{Prova 3 -- Turma 7}{26/10/2026}

% 01
%------------------------------------------------------------------------------%
\exercise{20\%}
Determine a posição relativa das retas.
\[
  r\colon\ \vetortri{x}{y}{z} = \vetortri{1}{0}{2} + t\vetortri{1}{1}{-1}
  \qquad\text{e}\qquad
  s\colon\ \vetortri{x}{y}{z} = \lambda\vetortri{2}{1}{0} + \vetortri{0}{1}{1}
\]

\begin{answer}
  Os vetores diretores $(1, 1, -1)$ e $(2, 1, 0)$ não são proporcionais.

  Para verificar se há interseção:
  \begin{align*}
    1+t & =2\lambda  \\
    t   & =1+\lambda \\
    2-t & =1
  \end{align*}
  A última equação dá $t=1$ e a segunda dá $\lambda=0$.

  A primeira exigiria $2=0$, impossível.

  Portanto, as retas são reversas.
\end{answer}

% 02
%------------------------------------------------------------------------------%
\exercise{20\%}
Determine a posição relativa das retas
\[
  r\colon (x, y, z) = (0, 1, 1) + t(1, -1, 2) \qquad
  s\colon (x, y, z) = (2, 0, 0) + \lambda(1, 1, 0)
\]

\begin{answer}
  Os vetores diretores $(1,-1,2)$ e $(1,1,0)$ não são proporcionais. Igualando as coordenadas:
  \begin{align*}
    t    & =2+\lambda \\
    1-t  & =\lambda   \\
    1+2t & =0
  \end{align*}
  A última dá $t=-1/2$ e a segunda dá $\lambda=3/2$

  A primeira exigiria $-1/2=7/2$, impossível

  As retas são reversas
\end{answer}

% 03
%------------------------------------------------------------------------------%
\exercise{20\%}
Determine o ângulo entre os vetores $\vec{v}=(1,0,1)$ e $\vec{w}=(0,1,1)$

\begin{answer}
  \[
    \vec{v}\cdot\vec{w}=1
    \qquad
    \norm{\vec{v}}=\norm{\vec{w}}=\sqrt{2}
  \]
  Logo $\cos\theta=1/2$ e $\theta=60^\circ$
\end{answer}

% 04
%------------------------------------------------------------------------------%
\exercise{20\%}
Verifique se $\vec{w}=(2, 5, 3)$
é combinação linear de $\vec{v}_1 = (1, 1, 0)$ e $\vec{v}_2=(0,1,1)$

\begin{answer}
  Procuramos $a$, $b$ tais que
  \[
    a(1,1,0)+b(0,1,1)=(2,5,3)
  \]
  Comparando coordenadas, $a=2$, $a+b=5$ e $b=3$, compatíveis.

  Portanto,
  \[
    \vec{w} = 2\vec{v}_1 + 3\vec{v}_2
  \]
\end{answer}

% 05
%------------------------------------------------------------------------------%
\exercise{20\%}
Obtenha, por Gram--Schmidt, uma base ortonormal que gere o mesmo espaço
que os vetores
$\vec{v}_1 = (1, 0)$ e
$\vec{v}_2 = (1, 1)$

\begin{answer}
  \[
    \vec{u}_1 = (1, 0)
    \qquad
    \vec{e}_1 = (1, 0)
  \]
  Em seguida,
  \begin{align*}
    \vec{u}_2 & =\vec{v}_2-\frac{\vec{v}_2\cdot\vec{u}_1}{\norm{\vec{u}_1}^{2}}\vec{u}_1 \\
              & =(1, 1) - (1, 0)                                                         \\
              & =(0, 1)
  \end{align*}
  Como $\norm{\vec{u}_2} = 1$, temos $\vec{e}_2 = (0, 1)$

  A base ortonormal é
  \[
    \{(1, 0), (0, 1)\}
  \]
\end{answer}

%------------------------------------------------------------------------------%
\end{document}
%------------------------------------------------------------------------------%

